% Section 2.8, from the summary table of lectures/L02/appendix/a_theory.md and from
% lectures/L02/README.md: the objectives, the evaluation questions and the next lesson.

\section{Sammanfattning}
\label{c2:sec:sammanfattning}

\begin{booktable}{@{}p{12.5em}L@{}}
  \thead{Regel} & \thead{Formel} \\
  \midrule
  Distributivlagen & $a(b + c) = ab + ac$ \\
  Kvadrat av summa & $(a + b)^2 = a^2 + 2ab + b^2$ \\
  Kvadrat av differens & $(a - b)^2 = a^2 - 2ab + b^2$ \\
  Konjugatregeln & $(a + b)(a - b) = a^2 - b^2$ \\
  Gemensam faktor & $ab + ac = a(b + c)$ \\
  Effekt & $P = UI = RI^2 = \frac{U^2}{R}$ \\
\end{booktable}

\needspace{6\baselineskip}
\subsection*{Det här ska du kunna}
Efter det här kapitlet ska du:
\begin{itemize}
  \item Förstå vad ett algebraiskt uttryck är och hur variabler används.
  \item Kunna förenkla algebraiska uttryck.
  \item Kunna multiplicera ut och faktorisera uttryck.
  \item Kunna använda reglerna för att förenkla och snabba upp beräkningar i elektriska kretsar.
\end{itemize}

\testadigsjalv
\begin{enumerate}
  \item Faktorisera $9R^2 - 4$.
  \item Förenkla $2(3I + 1) - (I - 5)$.
  \item Ett värmeelement med $R = 20\,\Omega$ matas först med $U_1 = 52\,\text{V}$ och sedan med
    $U_2 = 48\,\text{V}$. Använd konjugatregeln för att beräkna effektskillnaden
    \[
      P_1 - P_2 = \frac{U_1^2 - U_2^2}{R}
    \]
    utan miniräknare.
\end{enumerate}

\needspace{6\baselineskip}
\subsection*{Nästa kapitel}
\Kapref{3} handlar om linjära ekvationer och olikheter: hur de löses, hur lösningen till en olikhet
ritas på tallinjen och hur ekvationer med absolutbelopp hanteras.
